\section{Network and Tree Visualization}

This is a high-priority exam topic: graph-drawing algorithms appear under
``Visualization Algorithms'', and node-link\,$\leftrightarrow$\,matrix and
treemap\,$\leftrightarrow$\,node-link tree appear under ``Mapping data items'' /
``Alternatives''. The aesthetic-criteria list and the force-directed / layered /
morphing / foresighted layouts are recurring question targets.

\subsection{Graph basics}

\begin{keybox}[Definitions]
A \term{graph} $G=(V,E)$ is an abstract structure for modelling
\emph{relational} information: $V$ is a set of \term{nodes} (objects), $E$ a set
of \term{edges} (relations connecting nodes).
\begin{itemize}
  \item \term{Directed} vs.\ \term{undirected}: edges may have a direction. In a
  directed graph the \term{in-degree} of a node is the number of incoming edges,
  the \term{out-degree} the number of outgoing edges. The \term{degree} is the
  number of incident edges.
  \item \term{Weighted}: edges (or nodes) can carry \term{values} (edge weights),
  possibly of different types.
  \item Graphs may contain \term{cycles}; special families: planar, bipartite,
  hierarchical, clustered graphs, and \term{hypergraphs} (an edge connects more
  than two nodes).
  \item \term{Tree} = special case of a graph: \emph{no cycles}, one special
  \term{root} node, every non-root node has exactly one parent. Variants: free,
  binary, rooted, ordered trees.
  \item \term{Network} = graph + attributes on nodes/edges. A
  \term{multivariate network} attaches many attributes per element (primary data:
  directly measured; secondary data: derived/computed).
\end{itemize}
\end{keybox}

\paragraph{Data structures and the two visual representations.}
Internally a graph is stored as an \term{adjacency matrix} ($|V|\times|V|$, entry
$=1$ / weight if edge present) or an \term{adjacency list}. For \emph{display}
there are two fundamental mappings of the same data:

\begin{center}
\begin{minipage}{0.42\linewidth}\centering
\begin{tikzpicture}[every node/.style={font=\small}]
  \node[circle,draw,fill=tuwblue!15] (a) at (0,1.5) {1};
  \node[circle,draw,fill=tuwblue!15] (b) at (-0.9,0) {2};
  \node[circle,draw,fill=tuwblue!15] (c) at (0.9,0) {3};
  \draw (a)--(b); \draw (b)--(c);
\end{tikzpicture}\\[2pt]
{\small Node--link diagram}
\end{minipage}\hfill
\begin{minipage}{0.42\linewidth}\centering
\begin{tabular}{c|ccc}
  & 1 & 2 & 3 \\ \hline
1 & 0 & 1 & 0 \\
2 & 1 & 0 & 1 \\
3 & 0 & 1 & 0 \\
\end{tabular}\\[4pt]
{\small Adjacency matrix}
\end{minipage}
\end{center}

\subsection{Node-link vs.\ adjacency matrix (mapping \& alternatives)}

The two are the classic \term{alternative representations} of the \emph{same}
relational data, with complementary strengths.

\begin{center}
\begin{tabular}{p{0.30\linewidth} p{0.31\linewidth} p{0.31\linewidth}}
\toprule
 & \term{Node--link} & \term{Adjacency matrix} \\
\midrule
Edge $\to$ visual mark & line/curve between two node marks & filled cell
$(i,j)$ (colour/size $\to$ weight) \\
Path following / topology & \textbf{good} -- follow lines, see connectivity,
clusters by spatial proximity & \textbf{hard} -- must hop row$\to$column
mentally; multi-hop paths very tedious \\
Density / scalability & degrades on \textbf{dense} graphs (hairball, many edge
crossings, occlusion) & \textbf{good} on dense graphs: \emph{no edge crossings},
no overlap, every edge visible \\
Clusters / cliques & visible if layout separates them & a clique = solid block
on the diagonal; easy to spot with good row/column ordering \\
Space & compact for \emph{sparse} graphs & always $|V|^2$ cells \\
\bottomrule
\end{tabular}
\end{center}

\begin{keybox}[One-line rule of thumb]
\emph{Sparse} graph + need to \emph{follow paths / read topology} $\Rightarrow$
\term{node--link}. \emph{Dense} graph + want to spot \emph{clusters/cliques} and
avoid crossings $\Rightarrow$ \term{matrix}. Hybrid \term{NodeTrix} (Henry et
al.\ 2007) draws dense communities as small matrices and the sparse links
between them as node--link edges.
\end{keybox}

\subsection{Aesthetic criteria for graph drawing (the ``List'' topic)}

A good layout should be \emph{easy to read, understand and remember}. Layout
algorithms try to satisfy a set of \term{aesthetic criteria}; these are often
\emph{contradictory} and their optimisation is mostly \emph{NP-hard}, so most
graph-drawing algorithms are \term{heuristics}.

\begin{keybox}[Aesthetic criteria -- memorise the list]
\begin{itemize}[leftmargin=1.4em]
  \item Minimise \term{edge crossings} $\downarrow$
  \item Minimise \term{bends} of edges $\downarrow$ (max / uniform / total bends)
  \item Maximise \term{symmetry} $\uparrow$
  \item \term{Edge length} $\downarrow$: short, and \emph{uniform} edge lengths
  \item Even \term{node distribution}; \term{avoid node overlap}
  \item Minimise drawing \term{area} $\downarrow$
  \item Maximise \term{angular resolution} (``resolution'' $\uparrow$): angles
  between edges at a node should be large
  \item Show \term{clusters} / community structure
\end{itemize}
\end{keybox}

\begin{pitfall}[Criteria conflict]
The criteria \emph{trade off}. Classic exam example: a $K_4$ drawn as a
\emph{triangle with a centre node} has \emph{zero crossings} but no symmetry;
drawn as a \emph{square with both diagonals} it is fully \emph{symmetric} but has
one crossing. You cannot maximise everything at once -- say so explicitly.
\end{pitfall}

\subsection{Layout algorithm: Force-directed / spring embedder}

\begin{algobox}[Force-directed (spring embedder, Eades 1984)]
\textbf{Idea.} Model the graph as a \emph{physical system}: edges act as
\term{springs} (attractive force pulling adjacent nodes together, $\sim$ spring
toward a rest length), and \emph{all} pairs of nodes exert a \term{repulsive}
force (like charges / ``gravitation between nodes''). Iterate: compute net force
on each node, move it a little, repeat until the system reaches a
\term{stable, low-energy configuration} (energy minimisation).\\[2pt]
\textbf{Properties.} Produces \emph{symmetric, organic} layouts with
\emph{uniform edge lengths} -- exactly the two criteria Eades targeted. Works
well on small/medium sparse graphs; reveals clusters. Generalises to 3D.\\[2pt]
\textbf{Hyperparameters.} Spring constant / ideal (rest) edge length, repulsion
strength, step size or \term{temperature} (cooling schedule, as in
Fruchterman--Reingold), number of \term{iterations}, initial (often random)
placement.\\[2pt]
\textbf{Complexity.} Naive: $O(n^2)$ repulsion per iteration (all node pairs)
$\times$ iterations. \term{Barnes--Hut} approximation groups distant nodes into a
quad/oct-tree and treats each group as one pseudo-node $\Rightarrow$
$O(n\log n)$ per iteration.\\[2pt]
\textbf{Pitfalls.} \emph{Non-deterministic} (depends on random start);
gets stuck in \emph{local minima}; classic spring embedder has \emph{high
runtime} and can \emph{break down on very large graphs}. Improvements:
Kamada--Kawai 1989 (local minimal energies), Fruchterman--Reingold 1991, Davidson
\& Harel 1996 (simulated annealing).
\end{algobox}

\subsection{Layout algorithm: Layered / Sugiyama (hierarchies, flow)}

\begin{algobox}[Layered graph drawing (Sugiyama--Tagawa--Toda 1981)]
\textbf{Idea.} For \emph{directed} graphs / DAGs, hierarchies, and flow: arrange
nodes in horizontal \emph{layers} (rows) so all edges point downward. Classic
pipeline:
\begin{enumerate}[leftmargin=1.5em]
  \item \term{Cycle removal} -- extract an acyclic subgraph that contains all
  nodes (temporarily reverse a few edges).
  \item \term{Layer assignment} -- give each node an integer layer number; nodes
  arranged top-down in rows. (Long edges get \emph{dummy nodes} on intermediate
  layers.)
  \item \term{Crossing minimisation} -- choose the \emph{order} within each row to
  reduce edge crossings, typically layer-by-layer (e.g.\ \term{barycenter}
  heuristic: put a node at the average position of its neighbours in the adjacent
  layer).
  \item \term{Coordinate assignment} -- assign final $x$/$y$ to straighten edges
  and balance spacing.
\end{enumerate}
\textbf{Properties.} Emphasises \emph{direction / hierarchy / flow}; good for
DAGs, dependency and process diagrams. Variants: \emph{parallel layers}
(top-down rows) and \emph{radial layers} (concentric rings). \\[2pt]
\textbf{Hyperparameters.} Layer-assignment policy (longest-path, network-simplex),
crossing-minimisation heuristic and \#sweeps, node/layer spacing. \\[2pt]
\textbf{Complexity.} Even 2-layer crossing minimisation is \emph{NP-hard}
$\Rightarrow$ heuristics (barycenter, median). \\[2pt]
\textbf{Pitfalls.} Cycles must be handled first; dummy nodes inflate the graph;
poor for undirected / non-hierarchical graphs.
\end{algobox}

\subsection{Dynamic graphs and the mental map}

Networks increasingly \emph{change over time} (biochemical pathways gain new
reactions, social contacts appear/disappear). Visualisations of such dynamic
graphs must \term{preserve the ``Mental Map''} (Misue, Eades, Lai, Sugiyama
1995): the user should \emph{recognise old structures again} across time steps,
so nodes should not jump around between frames.

\begin{algobox}[Graph morphing]
\textbf{Idea.} Visualise the transition between two layouts with \emph{smooth
animation}: if a node changes position, compute an \emph{interpolated animation
path} between old and new position. Can be applied on top of \emph{any} 2D/3D
layout algorithm (e.g.\ GraphAEL uses force-based layouts).\\
\textbf{+} Looks good, eases tracking. \textbf{--} Nodes \emph{still change
position}; newly added or deleted nodes may not be correctly identified. This is
an \term{online} approach: lay out each step, then animate between consecutive
steps.
\end{algobox}

\begin{algobox}[Foresighted layout (Diehl, G\"org, Kerren 2001)]
\textbf{Idea.} If the \emph{whole sequence} of graphs is known (or can be
pre-computed) -- an \term{offline} setting -- build a \term{supergraph} (union of
all time steps) and compute \emph{one} layout from it. Position nodes at the
start so that they \emph{do not change their positions later} across the whole
sequence.\\
\textbf{+} \emph{Best mental-map preservation}; \emph{independent} of the chosen
layout algorithm. \textbf{--} The sequence is \emph{often unknown} in advance;
partly poor aesthetics (e.g.\ gaps at the beginning where future-only nodes are
reserved).
\end{algobox}

\begin{keybox}[Online vs.\ offline; animation vs.\ juxtaposition]
\term{Online}: layout per step as data arrives (morphing animates between them).
\term{Offline}: full sequence known $\Rightarrow$ foresighted layout for stable
positions. Time can be shown by \term{animation} (one view, moving) or by
\term{small multiples / timeline} (many static views side by side, e.g.\
\emph{matrix cubes}: stack adjacency matrices into a space-time cube).
\end{keybox}

\subsection{Tree layouts}

\begin{keybox}[Tree drawing families]
\begin{itemize}
  \item \term{Node--link trees}: \term{rectilinear} (classic top-down/left-right
  hierarchy with root on top) or \term{radial} (root in centre, depth $\to$
  radius).
  \item \term{Indented / outline} layout: nested indentation (like a file
  explorer); depth $\to$ horizontal indent.
  \item \term{Space-filling treemap}: \emph{recursive subdivision} of a rectangle.
  Hierarchy is encoded by \term{containment} (child rectangles nested inside the
  parent); a \emph{leaf}'s value is mapped to its \term{area}. No wasted
  whitespace. Algorithms: \term{slice-and-dice} (alternate cut direction per
  level -- simple but produces thin slivers) vs.\ \term{squarified} (keep cells
  close to square -- better aspect ratio, easier area comparison).
\end{itemize}
\end{keybox}

\begin{exambox}[Mapping data items: treemap $\leftrightarrow$ node-link tree]
\term{Treemap} and \term{rectilinear node-link tree} are two encodings of the
\emph{same hierarchy}. Node--link shows \emph{structure/depth} explicitly but
wastes space and scales poorly; treemap is \emph{space-filling} and great for
comparing \emph{leaf values by area} but makes the \emph{topology/depth} hard to
read. Containment (treemap) vs.\ connection (node--link) is the underlying
encoding difference.
\end{exambox}

\subsection{Biological networks (Kerren)}

Modern \term{systems biology} treats biological processes as integrated systems
of many interacting components; high-throughput methods (since the human genome
project) produce huge data, and \term{networks are the key concept} to structure
and combine it.

\begin{keybox}[Types of biological networks (a hierarchy)]
\begin{itemize}
  \item \term{Molecular graphs} -- atoms/bonds of a single molecule.
  \item \term{Metabolic networks / pathways} -- nodes are substances, edges are
  reactions, each single reaction catalysed by an \term{enzyme}; a \term{pathway}
  is a sequence of reactions converting one substance into another (databases:
  \emph{KEGG}; tools: PathViewer, BioPath, VANTED).
  \item \term{Protein--protein interaction (PPI) networks} -- nodes are proteins,
  edges are interactions (e.g.\ \emph{signal transduction}); fundamental for many
  functions and diseases (tool: \emph{Cytoscape}).
  \item \term{Gene regulatory networks} -- regulatory relations between genes.
\end{itemize}
Why it matters: diseases interpreted in a network context (infection = a foreign
network operating in ours; genetic defect = incorrect connectivity);
applications in drug design and metabolic engineering.
\end{keybox}

\begin{keybox}[Challenges \& integrative bioinformatics]
\textbf{Challenges:} \emph{scale} (huge networks, scalability/navigation);
\emph{integrating multiple heterogeneous data types} onto network elements
(multivariate primary/secondary attributes); respecting \emph{standard pathway-map
conventions} (e.g.\ KEGG layouts biologists expect).\\
\textbf{Integrative bioinformatics} combines network topology with the attached
multivariate data. State-of-the-art strategies (cf.\ the multivariate-network
slides):
\begin{itemize}
  \item \term{Integrated approaches} -- replace nodes by embedded diagrams
  (e.g.\ bar charts in pathway nodes).
  \item \term{Multiple coordinated views} -- network view linked with parallel
  coordinates, heatmaps, etc.\ (brushing-and-linking; e.g.\ Caleydo).
  \item \term{Semantic substrates} -- node position fixed by attributes in
  non-overlapping regions; edge slider controls clutter.
  \item \term{Attribute-driven topology} / hybrid (JauntyNets) -- user balances
  whether \emph{topology} or \emph{attributes} drive the force layout.
\end{itemize}
\end{keybox}

\subsection{Exam boxes}

\begin{exambox}[MCQ: node-link vs.\ adjacency matrix]
\emph{``For a large, very dense network where the analyst wants to find tightly
connected communities, which representation is preferable, and why?''} \\
Answer: the \term{adjacency matrix}. On dense graphs node--link becomes a
hairball with many crossings/occlusion, whereas a matrix has \emph{no edge
crossings}, shows \emph{every} edge, and -- with good row/column ordering --
makes cliques appear as solid diagonal blocks. Trade-off: \emph{path following}
across multiple hops is much harder in a matrix, so for sparse graphs / topology
tracing, prefer \term{node--link}.
\end{exambox}

\begin{exambox}[MCQ: force-directed vs.\ layered]
\emph{Which layout for an undirected social network where symmetry and uniform
edge length matter?} $\Rightarrow$ \term{force-directed}. \emph{Which for a
directed DAG / dependency or process flow?} $\Rightarrow$ \term{layered /
Sugiyama}. Distractor traps: force-directed is \emph{non-deterministic} and prone
to \emph{local minima} (true); Sugiyama's crossing minimisation is
\emph{NP-hard} hence heuristic (true); force-directed is naively $O(n^2)$ per
iteration, reducible to $O(n\log n)$ via \term{Barnes--Hut} (true).
\end{exambox}

\begin{exambox}[MCQ: morphing vs.\ foresighted layout]
Both aim to \term{preserve the mental map} of a \emph{dynamic} graph.
\term{Morphing} = \emph{online}, animate (interpolate) between per-step layouts;
nodes still move, can be applied to any layout, but added/removed nodes are hard
to follow. \term{Foresighted layout} = \emph{offline}, needs the full sequence,
computes a \emph{supergraph} so node positions \emph{stay stable} across all time
steps; best mental-map preservation but sequence often unknown and aesthetics can
suffer (gaps). Common wrong answer: claiming foresighted works online -- it does
\emph{not}; it requires the sequence in advance.
\end{exambox}
